% =====================================================================
\chapter{Soluciones}
\label{ap:soluciones}
% =====================================================================

Las soluciones de la parte~B se presentan con el esquema de los apuntes: datos, planteamiento, desarrollo, resultado y comprobación. Los valores intermedios se redondean a cuatro decimales; cuando el redondeo altera la cuarta cifra del resultado se indica el valor obtenido con todas las cifras. Todos los resultados numéricos se han comprobado con Python (NumPy y PyTorch). Cada solución termina con las secciones de los apuntes que conviene repasar si se ha fallado.

% ---------------------------------------------------------------------
\section{Clave de la parte A}
\label{sec:clave}

\begin{center}
\small
\renewcommand{\arraystretch}{1.12}
\begin{tabularx}{\textwidth}{@{}c c >{\raggedright\arraybackslash}X l@{}}
\toprule
Nº & Resp. & Justificación & Apuntes \\
\midrule
1 & (b) & $f''=4=L$; el método converge si y solo si $\abs{1-4\eta}<1$, es decir, $0<\eta<2/L=0{,}5$. & M07, §2.3 \\
2 & (c) & La acumulación exponencial de un valor constante tiende a $g/(1-\beta)=g/0{,}05=20\,g$. & M07, §4.1 \\
3 & (b) & $\Var[\hat g]=\sigma^{2}/B$; la desviación típica escala como $1/\sqrt B$ y $\sqrt4=2$. & M07, §3.1 \\
4 & (a) & El decaimiento desacoplado resta $\eta\lambda\theta=10^{-3}\cdot 0{,}01\cdot 4=4\cdot 10^{-5}$, sin dividir por $\sqrt{\hat v_t}$. & M07, §6.2 \\
5 & (b) & Al principio los momentos de Adam se basan en muy pocos gradientes, que además son grandes e irregulares. & M07, §6.3 \\
6 & (c) & Anular el gradiente da las ecuaciones normales $X\T X\theta=X\T y$. & M08, §2.2 \\
7 & (a) & $\sigma'$ se cancela con la derivada del logaritmo y queda $(p-y)\,x$. & M08, §3.4 \\
8 & (c) & Con logits iguales la pérdida es $\ln K=\ln 28=3{,}3322$ nats. & M08, §4.5 \\
9 & (b) & Error bajo en entrenamiento y alto en validación: el modelo memoriza y no generaliza. & M08, §5.3 \\
10 & (b) & \texttt{nn.Linear(n, m)} guarda \texttt{weight} con forma $(m,n)$, traspuesta de $W$. & M10, §6.2 \\
11 & (b) & $s=\sqrt{2/(100+200)}=\sqrt{1/150}=0{,}0816$. & M09, §4.3 \\
12 & (b) & $\GELU(-1)=-0{,}1587<0$, mientras que $\ReLU(-1)=0$. & M09, §2.3 \\
13 & (b) & Tres matrices de $4096\times 11\,008$: $3\cdot 45\,088\,768=135\,266\,304$. & M09, §2.4 \\
14 & (b) & Cada pasada inversa da una fila del jacobiano; hay $m=3$ salidas. & M10, §3.1 \\
15 & (c) & $(2+\varepsilon)^{3}=8+3\cdot 2^{2}\varepsilon=8+12\varepsilon$. & M10, §2.2 \\
16 & (b) & $\RMS(10x)=10\RMS(x)$ y el cociente no cambia (invariancia de escala). & M11, §3.2 \\
17 & (b) & Se anula con probabilidad $0{,}25$ o se divide por $0{,}75$: $4/0{,}75=5{,}3333$. & M11, §4.1 \\
18 & (c) & $67\,600$ supera el máximo de fp16, $65\,504$, y se redondea a \Inf. (En bf16 daría $67\,584$.) & M11, §5.1 \\
19 & (a) & $\ln 65\,504=11{,}09$. (El valor $88{,}72$ es el umbral de fp32.) & M11, §6.1 \\
20 & (b) & Salto brusco, sin crecimiento previo y sin fp16: operación sin proteger. & M11, §6.4 \\
\bottomrule
\end{tabularx}
\end{center}

% ---------------------------------------------------------------------
\section{Soluciones de la parte B}
\label{sec:sol-problemas}

% ===================================================================== P1
\begin{solucion}{eje:mlp}
\paragraph{Apartado (a).}
\datos $X$, $y=(0,2)$, $W_1$, $b_1=(0;\,-0{,}5)$, $W_2$, $b_2=0$.
\planteamiento pasada hacia delante fila a fila (M09, §5.5). La pérdida de cada ejemplo se calcula como $\ell_i=\LSE\bigl(z^{(2)}_i\bigr)-z^{(2)}_{i,y_i}$, que es más precisa que $-\ln P_{i,y_i}$ con $P$ redondeada (M11, §6.3).
\pasos
\begin{enumerate}[label=\arabic*.]
  \item Fila~1 de $XW_1$: $(2\cdot 1+1\cdot(-1);\allowbreak\ 2\cdot 0{,}5+1\cdot 1)=(1;\allowbreak\ 2)$; con $b_1$, $(1;\allowbreak\ 1{,}5)$.
  \item Fila~2 de $XW_1$: $(1+1;\allowbreak\ 0{,}5-1)=(2;\allowbreak\ -0{,}5)$; con $b_1$, $(2;\allowbreak\ -1)$.
  \item La ReLU anula el $-1$. $HW_2$: fila~1, $(1;\allowbreak\ 1{,}5;\allowbreak\ -1+1{,}5)=(1;\allowbreak\ 1{,}5;\allowbreak\ 0{,}5)$; fila~2, $(2;\allowbreak\ 0;\allowbreak\ -2)$.
  \[
  Z^{(1)}=\begin{pmatrix}1&1{,}5\\2&-1\end{pmatrix},\qquad H=\begin{pmatrix}1&1{,}5\\2&0\end{pmatrix},\qquad Z^{(2)}=\begin{pmatrix}1&1{,}5&0{,}5\\2&0&-2\end{pmatrix}.
  \]
  \item Fila~1: $e^{z}=(2{,}7183;\allowbreak\ 4{,}4817;\allowbreak\ 1{,}6487)$, suma $8{,}8487$; $P_1=(0{,}3072;\allowbreak\ 0{,}5065;\allowbreak\ 0{,}1863)$.
  \item Fila~2: $e^{z}=(7{,}3891;\allowbreak\ 1;\allowbreak\ 0{,}1353)$, suma $8{,}5244$; $P_2=(0{,}8668;\allowbreak\ 0{,}1173;\allowbreak\ 0{,}0159)$.
  \item $\ell_1=\ln 8{,}8487-1=2{,}1803-1=1{,}1803$; $\ell_2=\ln 8{,}5244-(-2)=2{,}1429+2=4{,}1429$.
  \item $L=(1{,}1803+4{,}1429)/2=2{,}6616$.
\end{enumerate}
\resultado $P=\begin{pmatrix}0{,}3072&0{,}5065&0{,}1863\\0{,}8668&0{,}1173&0{,}0159\end{pmatrix}$, $\ell_1=1{,}1803$, $\ell_2=4{,}1429$, $L=2{,}6616$.
\comprobacion cada fila de $P$ suma~1. Con $P$ redondeada, $-\ln 0{,}0159=4{,}1414$: por eso se usa $\LSE-z$. El segundo ejemplo, de la clase~2, recibe solo un $1{,}6\,\%$ de probabilidad.

\paragraph{Apartado (b).}
\datos $p=(0{,}3072;\allowbreak\ 0{,}5065;\allowbreak\ 0{,}1863)$, $y_1=0$, $\ell_1=-\ln p_0$.
\planteamiento $J_{ij}=\partial p_i/\partial z_j=p_i(\uno[i=j]-p_j)$ y, por la regla de la cadena con vectores fila, $\partial\ell_1/\partial z=(\partial\ell_1/\partial p)\,J$ (M09, §5.2).
\pasos
\begin{enumerate}[label=\arabic*.]
  \item Diagonal: $p_0(1-p_0)=0{,}2128$, $p_1(1-p_1)=0{,}2500$, $p_2(1-p_2)=0{,}1516$. Fuera de la diagonal: $-p_0p_1=-0{,}1556$, $-p_0p_2=-0{,}0572$, $-p_1p_2=-0{,}0944$.
  \[
  J=\begin{pmatrix}0{,}2128&-0{,}1556&-0{,}0572\\-0{,}1556&0{,}2500&-0{,}0944\\-0{,}0572&-0{,}0944&0{,}1516\end{pmatrix},\qquad
  \frac{\partial\ell_1}{\partial p}=\Bigl(-\frac{1}{p_0};\allowbreak\ 0;\allowbreak\ 0\Bigr)=(-3{,}2553;\allowbreak\ 0;\allowbreak\ 0).
  \]
  \item Solo cuenta la primera fila de $J$:
  \[
  -3{,}2553\cdot(0{,}2128;\allowbreak\ -0{,}1556;\allowbreak\ -0{,}0572)=(-0{,}6927;\allowbreak\ 0{,}5065;\allowbreak\ 0{,}1862),
  \]
  que coincide con $p-e_0=(-0{,}6928;\allowbreak\ 0{,}5065;\allowbreak\ 0{,}1863)$ salvo redondeo en la cuarta cifra.
  \item Caso general: si $y=e_k$, $\mathrm{CE}=-\ln\hat y_k=-\theta_k+\LSE(\theta)$. Como $\partial\LSE/\partial\theta_j=e^{\theta_j}/\sum_m e^{\theta_m}=\hat y_j$, se obtiene $\partial\mathrm{CE}/\partial\theta_j=\hat y_j-\uno[j=k]$, es decir, $\partial\mathrm{CE}/\partial\theta=\hat y-y$.
  \item Factor $1/2$: $L=(\ell_1+\ell_2)/2$ y $\ell_i$ solo depende de la fila~$i$ de $Z^{(2)}$, luego $\partial L/\partial Z^{(2)}_i=\tfrac12(P_i-e_{y_i})$. Apilando las filas, $\Delta^{(2)}=(P-Y)/2$.
\end{enumerate}
\resultado $(\partial\ell_1/\partial p)\,J=p-e_0\approx(-0{,}6928;\allowbreak\ 0{,}5065;\allowbreak\ 0{,}1863)$; el factor $1/2$ procede de promediar dos ejemplos.
\comprobacion las filas de $J$ suman~0 ($0{,}2128-0{,}1556-0{,}0572=0$), por lo que el gradiente también suma~0.

\paragraph{Apartado (c).}
\datos $P$ y $H$ de (a); $Y=\begin{pmatrix}1&0&0\\0&0&1\end{pmatrix}$.
\planteamiento $\Delta^{(2)}=(P-Y)/2$; $\nabla W_2=H\T\Delta^{(2)}$; $\nabla b_2=\uno\T\Delta^{(2)}$ (M09, §5.3).
\pasos
\begin{enumerate}[label=\arabic*.]
  \item $P-Y=\begin{pmatrix}-0{,}6928&0{,}5065&0{,}1863\\0{,}8668&0{,}1173&-0{,}9841\end{pmatrix}$, luego $\Delta^{(2)}=\begin{pmatrix}-0{,}3464&0{,}2532&0{,}0932\\0{,}4334&0{,}0587&-0{,}4921\end{pmatrix}$.
  \item Fila~1 de $\nabla W_2$: $1\cdot\Delta^{(2)}_1+2\cdot\Delta^{(2)}_2=(0{,}5204;\allowbreak\ 0{,}3706;\allowbreak\ -0{,}8910)$.
  \item Fila~2 de $\nabla W_2$: $1{,}5\cdot\Delta^{(2)}_1+0\cdot\Delta^{(2)}_2=(-0{,}5196;\allowbreak\ 0{,}3799;\allowbreak\ 0{,}1397)$.
  \item $\nabla b_2=(-0{,}3464+0{,}4334;\allowbreak\ 0{,}2532+0{,}0587;\allowbreak\ 0{,}0932-0{,}4921)=(0{,}0870;\allowbreak\ 0{,}3119;\allowbreak\ -0{,}3989)$.
\end{enumerate}
\resultado $\nabla W_2=\begin{pmatrix}0{,}5204&0{,}3706&-0{,}8910\\-0{,}5196&0{,}3799&0{,}1397\end{pmatrix}$, $\nabla b_2=(0{,}0870;\allowbreak\ 0{,}3119;\allowbreak\ -0{,}3989)$.
\comprobacion $H\T$ ($2\times2$) por $\Delta^{(2)}$ ($2\times3$) da $2\times3$, como $W_2$. Las filas de $\nabla W_2$ y $\nabla b_2$ suman~0 porque las de $\Delta^{(2)}$ también.

\paragraph{Apartado (d).}
\datos $\Delta^{(2)}$, $W_2$, $Z^{(1)}$ y $X$.
\planteamiento la ReLU deja pasar el gradiente solo donde $Z^{(1)}>0$ (M09, §5.4); $\nabla W_1=X\T\Delta^{(1)}=\sum_i x_i\T\Delta^{(1)}_i$.
\pasos
\begin{enumerate}[label=\arabic*.]
  \item $\nabla H=\Delta^{(2)}W_2\T$: fila~1, $(-0{,}3464-0{,}0932;\allowbreak\ 0{,}2532+0{,}0932)=(-0{,}4396;\allowbreak\ 0{,}3464)$; fila~2, $(0{,}4334+0{,}4921;\allowbreak\ 0{,}0587-0{,}4921)=(0{,}9255;\allowbreak\ -0{,}4334)$.
  \item Máscara $\uno[Z^{(1)}>0]=\begin{pmatrix}1&1\\1&0\end{pmatrix}$: está apagada la neurona~2 en el ejemplo~2 ($Z^{(1)}_{22}=-1$). Luego $\Delta^{(1)}=\begin{pmatrix}-0{,}4396&0{,}3464\\0{,}9255&0\end{pmatrix}$.
  \item Fila~1 de $\nabla W_1$: $2\cdot(-0{,}4396;\allowbreak\ 0{,}3464)+1\cdot(0{,}9255;\allowbreak\ 0)=(0{,}0463;\allowbreak\ 0{,}6928)$. Fila~2: $1\cdot(-0{,}4396;\allowbreak\ 0{,}3464)-1\cdot(0{,}9255;\allowbreak\ 0)=(-1{,}3650;\allowbreak\ 0{,}3464)$.
  \item $\nabla b_1=(-0{,}4396+0{,}9255;\allowbreak\ 0{,}3464+0)=(0{,}4859;\allowbreak\ 0{,}3464)$.
  \item Columna~2 de $\nabla W_1$: $x_1\T\Delta^{(1)}_{12}+x_2\T\Delta^{(1)}_{22}=(2;\,1)\T\cdot 0{,}3464+(1;\,-1)\T\cdot 0=(0{,}6928;\allowbreak\ 0{,}3464)\T$: solo contribuye el primer ejemplo.
\end{enumerate}
\resultado $\nabla W_1=\begin{pmatrix}0{,}0463&0{,}6928\\-1{,}3650&0{,}3464\end{pmatrix}$, $\nabla b_1=(0{,}4859;\allowbreak\ 0{,}3464)$; la unidad apagada es la neurona~2 en el ejemplo~2.
\comprobacion las diferencias centradas con paso $10^{-6}$ (M09, §5.6) reproducen $\nabla W_1$ con cuatro decimales. La neurona~2 no está \emph{muerta}: lo estaría si se apagara con todos los ejemplos (M09, §2.2).

\paragraph{Apartado (e).}
\datos $\Delta^{(1)}$ de (d) y $W_1$.
\planteamiento $Z^{(1)}=XW_1+b_1$ es afín en $X$; por la regla de una capa afín (M09, §5.3), $\nabla_X L=\Delta^{(1)}W_1\T$.
\pasos
\begin{enumerate}[label=\arabic*.]
  \item $W_1\T=\begin{pmatrix}1&-1\\0{,}5&1\end{pmatrix}$.
  \item Fila~1: $(-0{,}4396+0{,}5\cdot 0{,}3464;\allowbreak\ 0{,}4396+0{,}3464)=(-0{,}2664;\allowbreak\ 0{,}7860)$.
  \item Fila~2: $(0{,}9255+0;\allowbreak\ -0{,}9255+0)=(0{,}9255;\allowbreak\ -0{,}9255)$.
\end{enumerate}
\resultado $\nabla_X L=\Delta^{(1)}W_1\T=\begin{pmatrix}-0{,}2664&0{,}7860\\0{,}9255&-0{,}9255\end{pmatrix}$. Las mayores sensibilidades, en valor absoluto $0{,}9255$, corresponden a las dos componentes del segundo ejemplo, el peor clasificado; en el primero domina la segunda componente.
\comprobacion las diferencias centradas dan el mismo resultado. Es la expresión $\partial J/\partial x=\delta^{(1)}W_1\T$ de la fuente, con $\delta^{(1)}=\delta^{(2)}W_2\T\odot\sigma'(z^{(1)})$ y la sigmoide sustituida por la ReLU.

\paragraph{Apartado (f).}
\datos gradientes de (c) y (d); $\eta=0{,}5$.
\planteamiento $\theta\leftarrow\theta-\eta\nabla_\theta L$ en los cuatro tensores y nueva pasada hacia delante.
\pasos
\begin{enumerate}[label=\arabic*.]
  \item $W_1=\begin{pmatrix}0{,}9768&0{,}1536\\-0{,}3175&0{,}8268\end{pmatrix}$, $b_1=(-0{,}2430;\allowbreak\ -0{,}6732)$.
  \item $W_2=\begin{pmatrix}0{,}7398&-0{,}1853&-0{,}5545\\0{,}2598&0{,}8101&0{,}9301\end{pmatrix}$, $b_2=(-0{,}0435;\allowbreak\ -0{,}1559;\allowbreak\ 0{,}1995)$.
  \item $Z^{(1)}=\begin{pmatrix}1{,}3932&0{,}4608\\1{,}0514&-1{,}3464\end{pmatrix}$ y $Z^{(2)}=\begin{pmatrix}1{,}1069&-0{,}0408&-0{,}1445\\0{,}7343&-0{,}3507&-0{,}3835\end{pmatrix}$.
  \item Fila~1: $e^{z}=(3{,}0250;\allowbreak\ 0{,}9600;\allowbreak\ 0{,}8654)$, suma $4{,}8505$; $P_1=(0{,}6237;\allowbreak\ 0{,}1979;\allowbreak\ 0{,}1784)$; $\ell_1=-\ln 0{,}6237=0{,}4722$.
  \item Fila~2: $e^{z}=(2{,}0840;\allowbreak\ 0{,}7042;\allowbreak\ 0{,}6814)$, suma $3{,}4696$; $P_2=(0{,}6006;\allowbreak\ 0{,}2030;\allowbreak\ 0{,}1964)$; $\ell_2=\ln 3{,}4696+0{,}3835=1{,}6276$.
  \item $L=(0{,}4722+1{,}6276)/2=1{,}0499$.
\end{enumerate}
\resultado $\ell_1=0{,}4722$ (antes $1{,}1803$), $\ell_2=1{,}6276$ (antes $4{,}1429$), $L=1{,}0499$ (antes $2{,}6616$). Mejoran los dos ejemplos.
\comprobacion con los parámetros sin redondear, PyTorch da $L=1{,}0499$. La bajada es grande porque $\eta=0{,}5$ es un paso generoso para una red tan pequeña.

\paragraph{Apartado (g).}
\datos red $D_x\to H_o\to D_y$ con sesgos; aquí $D_x=2$, $H_o=2$, $D_y=3$.
\planteamiento cada capa afín $n\to m$ tiene $nm$ pesos y $m$ sesgos; si la pérdida es la suma, $L_{\text{suma}}=B\,L_{\text{media}}$.
\pasos
\begin{enumerate}[label=\arabic*.]
  \item Total: $D_xH_o+H_o+H_oD_y+D_y=(D_x+1)H_o+(H_o+1)D_y$.
  \item Aquí: $(2+1)\cdot 2+(2+1)\cdot 3=6+9=15$.
  \item Con la suma, $\Delta^{(2)}=P-Y$ sin el factor $1/2$ y, por linealidad de la retropropagación, todos los gradientes se multiplican por $B=2$.
  \item Con la suma, la escala del gradiente crece con $B$ y habría que reajustar $\eta$ al cambiar el tamaño del lote; con la media, el gradiente de minibatch es un estimador insesgado de $\nabla J$ cuya escala no depende de $B$ (M07, §3.1).
\end{enumerate}
\resultado \textbf{$(D_x+1)H_o+(H_o+1)D_y$ parámetros, 15 en esta red; con la suma, todos los gradientes se multiplican por $B$; la media hace que $\eta$ no dependa de $B$.}
\comprobacion con la suma, $\nabla b_2=2\cdot(0{,}0870;\allowbreak\ 0{,}3119;\allowbreak\ -0{,}3989)=(0{,}1740;\allowbreak\ 0{,}6238;\allowbreak\ -0{,}7978)$.

\paragraph{Criterios de corrección.} (a) 1~p $Z^{(1)}$ y $H$, 1~p $Z^{(2)}$, 1~p $P$ y pérdidas. (b) 1~p jacobiano y producto, 1~p demostración general, 1~p factor $1/2$. (c) 1~p por cada resultado. (d) 1~p $\Delta^{(1)}$ y unidad apagada, 1~p $\nabla W_1$ y $\nabla b_1$, 1~p columna~2. (e) 1~p fórmula, 1~p valores. (f) 1~p parámetros, 1~p pérdidas. (g) 1~p conteo, 1~p suma frente a media.

\textit{Repasar:} M09, §5.2--5.5; M08, §4.3; M09, §5.6; M07, §3.1.
\end{solucion}

% ===================================================================== P2
\begin{solucion}{eje:opt}
\paragraph{Apartado (a).}
\datos $J(w)=w\T Aw+b\T w$, $A=\diag(2,\,1/2)$, $b=(-4,\,-1)\T$, $w^{(0)}=0$.
\planteamiento $A$ es simétrica, luego $\nabla J=2Aw+b$ y la hessiana es $2A$. En una cuadrática el descenso de gradiente es lineal en el error y se analiza por valores propios (M07, §2.3; \cite[§9.3]{bv}, texto del curso EE364a~\cite{ee364a}).
\pasos
\begin{enumerate}[label=\arabic*.]
  \item $\nabla J(w)=2Aw+b=(4w_1-4;\allowbreak\ w_2-1)\T$. Se anula en $w^{*}=-\tfrac12A^{-1}b=(1,\,1)\T$.
  \item Restando $w^{*}$ a la iteración y usando $2Aw^{*}+b=0$: $e^{(k+1)}=e^{(k)}-\eta(2Aw^{(k)}+b)=e^{(k)}-\eta\,2A\bigl(w^{(k)}-w^{*}\bigr)=(I-2\eta A)\,e^{(k)}$.
  \item $I-2\eta A=\diag(1-4\eta,\ 1-\eta)$. El error tiende a~0 desde todo $e^{(0)}$ si y solo si $\abs{1-4\eta}<1$ y $\abs{1-\eta}<1$, es decir, $0<\eta<0{,}5$. Equivale a la condición de la fuente, $0<\lambda_i(A)<1/\eta$.
  \item Hessiana $2A=\diag(4,1)$: $L=4$, $\mu=1$, $\kappa=L/\mu=4$.
  \item $\eta^{*}=2/(\mu+L)=2/5=0{,}4$; factor $(\kappa-1)/(\kappa+1)=3/5=0{,}6$. Los factores por coordenada son $1-0{,}4\cdot 4=-0{,}6$ y $1-0{,}4=0{,}6$.
  \item $e^{(0)}=(-1,\,-1)\T$; $e^{(1)}=(0{,}6;\allowbreak\ -0{,}6)\T$, luego $w^{(1)}=(1{,}6;\allowbreak\ 0{,}4)\T$; $e^{(2)}=(-0{,}36;\allowbreak\ -0{,}36)\T$, luego $w^{(2)}=(0{,}64;\allowbreak\ 0{,}64)\T$.
\end{enumerate}
\resultado $w^{*}=(1,1)\T$; \textbf{convergencia si y solo si $0<\eta<0{,}5$}; $\kappa=4$; $\eta^{*}=0{,}4$ \textbf{con factor} $0{,}6$; $w^{(1)}=(1{,}6;\allowbreak\ 0{,}4)\T$, $w^{(2)}=(0{,}64;\allowbreak\ 0{,}64)\T$.
\comprobacion directamente, $\nabla J(w^{(1)})=(2{,}4;\allowbreak\ -0{,}6)\T$ y $w^{(1)}-0{,}4\cdot(2{,}4;\allowbreak\ -0{,}6)\T=(0{,}64;\allowbreak\ 0{,}64)\T$. La primera coordenada del error cambia de signo en cada paso y la segunda no; ambas se reducen en un factor $0{,}6$.

\paragraph{Apartado (b).}
\datos $\theta_0=2$, $\eta=0{,}1$, $\lambda=0{,}1$, $\beta_1=0{,}9$, $\beta_2=0{,}999$, $\varepsilon=0$, $g_1=4$, $g_2=1$.
\planteamiento Adam (M07, §5.3) con el decaimiento desacoplado de AdamW (M07, §6.2): $\theta_t=\theta_{t-1}-\eta\bigl(\hat m_t/\sqrt{\hat v_t}+\lambda\theta_{t-1}\bigr)$.
\pasos
\begin{enumerate}[label=\arabic*.]
  \item $t=1$: $m_1=0{,}1\cdot 4=0{,}4$; $v_1=0{,}001\cdot 16=0{,}016$; $\hat m_1=0{,}4/0{,}1=4$; $\hat v_1=0{,}016/0{,}001=16$.
  \item $\hat m_1/\sqrt{\hat v_1}=4/4=1$; $\lambda\theta_0=0{,}2$; paso $0{,}1\cdot 1{,}2=0{,}12$; $\theta_1=1{,}88$.
  \item $t=2$: $m_2=0{,}36+0{,}1=0{,}46$; $v_2=0{,}015984+0{,}001=0{,}016984$; $\hat m_2=0{,}46/0{,}19=2{,}4211$; $\hat v_2=0{,}016984/0{,}001999=8{,}4962$.
  \item $\sqrt{8{,}4962}=2{,}9148$; $\hat m_2/\sqrt{\hat v_2}=0{,}8306$; $\lambda\theta_1=0{,}188$; paso $0{,}1\cdot 1{,}0186=0{,}1019$; $\theta_2=1{,}7781$.
\end{enumerate}
\begin{center}\small
\begin{tabular}{@{}cccccccc@{}}
\toprule
$t$ & $g_t$ & $m_t$ & $v_t$ & $\hat m_t$ & $\hat v_t$ & $\hat m_t/\sqrt{\hat v_t}$ & $\theta_t$ \\
\midrule
1 & 4 & $0{,}4$ & $0{,}016$ & $4$ & $16$ & $1$ & $1{,}88$ \\
2 & 1 & $0{,}46$ & $0{,}016984$ & $2{,}4211$ & $8{,}4962$ & $0{,}8306$ & $1{,}7781$ \\
\bottomrule
\end{tabular}
\end{center}
\resultado $\theta_1=1{,}88$, $\theta_2=1{,}7781$.
\comprobacion en $t=1$, $\hat m_1/\sqrt{\hat v_1}=g_1/\abs{g_1}=1$ (M07, §5.3). \texttt{torch.optim.AdamW} con estos datos da $1{,}8800$ y $1{,}7781$.

\paragraph{Apartado (c).}
\datos $\eta_{\max}=10^{-3}$, $\eta_{\min}=10^{-4}$, $W=500$, $T=5500$.
\planteamiento $\eta_t=\eta_{\max}t/W$ si $t<W$, y $\eta_t=\eta_{\min}+\tfrac12(\eta_{\max}-\eta_{\min})\bigl(1+\cos(\pi(t-W)/(T-W))\bigr)$ si $W\le t\le T$ (M07, §6.3).
\pasos
\begin{enumerate}[label=\arabic*.]
  \item $t=250$: $10^{-3}\cdot 250/500=5\cdot 10^{-4}$.
  \item $\tfrac12(\eta_{\max}-\eta_{\min})=4{,}5\cdot 10^{-4}$ y $T-W=5000$.
  \item $t=3000$: cociente $0{,}5$, $\cos(\pi/2)=0$: $10^{-4}+4{,}5\cdot 10^{-4}=5{,}5\cdot 10^{-4}$.
  \item $t=4250$: cociente $0{,}75$, $\cos(3\pi/4)=-0{,}7071$: $10^{-4}+4{,}5\cdot 10^{-4}\cdot 0{,}2929=2{,}318\cdot 10^{-4}$.
  \item $t=5500$: $\cos\pi=-1$: $10^{-4}$.
\end{enumerate}
\resultado $5\cdot 10^{-4}$, $5{,}5\cdot 10^{-4}$, $2{,}318\cdot 10^{-4}$ \textbf{y} $10^{-4}$.
\comprobacion en el punto medio del tramo coseno, $\eta=(\eta_{\max}+\eta_{\min})/2=5{,}5\cdot 10^{-4}$, como establece la proposición de M07, §6.3.

\paragraph{Criterios de corrección.} (a) 1~p gradiente y $w^{*}$, 1~p recurrencia, 1~p condición y $\kappa$, 1~p $\eta^{*}$ y factor, 1~p iterados. (b) 1,5~p por fila. (c) 0,5~p por valor.

\textit{Repasar:} M07, §2.2--2.3; M07, §5.3; M07, §6.2--6.3.
\end{solucion}

% ===================================================================== P3
\begin{solucion}{eje:sup}
\paragraph{Apartado (a).}
\datos $x=(1,2,3)$, $y=(1,2,6)$; $X$ tiene filas $(1,\,x_i)$.
\planteamiento ecuaciones normales $X\T X\theta=X\T y$ (M08, §2.2) y $\hat\sigma^{2}=\SSE/n$ (M08, §2.5; \cite[§1.3]{cs229}).
\pasos
\begin{enumerate}[label=\arabic*.]
  \item $X\T X=\begin{pmatrix}3&6\\6&14\end{pmatrix}$ y $X\T y=(9;\allowbreak\ 23)\T$.
  \item $\det=42-36=6$; $\hat\theta=\tfrac16\begin{pmatrix}14&-6\\-6&3\end{pmatrix}\begin{pmatrix}9\\23\end{pmatrix}=\tfrac16\begin{pmatrix}-12\\15\end{pmatrix}=\begin{pmatrix}-2\\2{,}5\end{pmatrix}$.
  \item Predicciones $(0{,}5;\allowbreak\ 3;\allowbreak\ 5{,}5)$; residuos $y-\hat y=(0{,}5;\allowbreak\ -1;\allowbreak\ 0{,}5)$.
  \item $\SSE=0{,}25+1+0{,}25=1{,}5$; $\hat\sigma^{2}=1{,}5/3=0{,}5$.
\end{enumerate}
\resultado $\hat\theta=(-2;\allowbreak\ 2{,}5)$, \textbf{residuos} $(0{,}5;\allowbreak\ -1;\allowbreak\ 0{,}5)$, $\SSE=1{,}5$, $\hat\sigma^{2}=0{,}5$.
\comprobacion el vector de residuos es ortogonal a las columnas de $X$ (M08, §2.3): suma~0 y $1\cdot 0{,}5+2\cdot(-1)+3\cdot 0{,}5=0$.

\paragraph{Apartado (b).}
\datos $P(y=1\mid x)=\sigma(x-\alpha)$; $J(\alpha)=-y\ln s-(1-y)\ln(1-s)$ con $s=\sigma(x-\alpha)$.
\planteamiento regla de la cadena con $\mathrm{d}s/\mathrm{d}\alpha=-s(1-s)$ (M08, §3.1) y convexidad por el signo de $J''$ (M08, §3.4).
\pasos
\begin{enumerate}[label=\arabic*.]
  \item $J'(\alpha)=-y\,\dfrac{-s(1-s)}{s}-(1-y)\,\dfrac{s(1-s)}{1-s}=y(1-s)-(1-y)s=y-s$.
  \item SGD: $\alpha\leftarrow\alpha-\epsilon J'(\alpha)=\alpha+\epsilon\bigl(\sigma(x-\alpha)-y\bigr)$. Con $(x,y)=(3,1)$ y $\alpha=1$: $s=\sigma(2)=0{,}8808$ y $\alpha\leftarrow 1+0{,}5\,(0{,}8808-1)=0{,}9404$.
  \item $J''(\alpha)=-\mathrm{d}s/\mathrm{d}\alpha=s(1-s)>0$ porque $s\in(0,1)$: $J$ es estrictamente convexa en $\R$.
  \item Con $(0,0)$ y $(2,1)$: la derivada de la suma es $\bigl(0-\sigma(-\alpha)\bigr)+\bigl(1-\sigma(2-\alpha)\bigr)$. En $\alpha=1$ vale $1-\sigma(-1)-\sigma(1)=1-0{,}2689-0{,}7311=0$, porque $\sigma(-t)=1-\sigma(t)$. Como la suma es estrictamente convexa, $\alpha^{*}=1$ es su único mínimo.
  \item Si todos los ejemplos son de la clase~1, por ejemplo $\{(0,1),(1,1)\}$, la derivada de la suma es $\sum_i\bigl(1-\sigma(x_i-\alpha)\bigr)>0$: la pérdida es estrictamente creciente en $\alpha$, tiende a~0 cuando $\alpha\to-\infty$ y ese ínfimo no se alcanza. Lo mismo ocurre, con la desigualdad contraria, si todos son de la clase~0.
\end{enumerate}
\resultado $J'(\alpha)=y-\sigma(x-\alpha)$; \textbf{tras un paso,} $\alpha=0{,}9404$; $J''=s(1-s)>0$ \textbf{(convexa)}; $\alpha^{*}=1$; \textbf{sin mínimo si todos los ejemplos son de la misma clase.}
\comprobacion el paso desplaza la frontera hacia la izquierda, de modo que $x=3$ queda más claramente en la clase~1. El mínimo $\alpha^{*}=1$ es el punto medio entre $0$ y $2$, como exige la simetría del problema. La situación del paso~5 es la versión unidimensional del caso de datos separables de M08, §3.4: la pérdida tiene ínfimo pero no mínimo, y el SGD desplaza $\alpha$ sin límite.

\paragraph{Apartado (c).}
\datos $z=(1,\,0,\,-1,\,2)$, clase correcta $3$.
\planteamiento $p=\softmax(z)$, $\ell=\LSE(z)-z_3$ y $\nabla_z\ell=p-e_3$ (M08, §4.3).
\pasos
\begin{enumerate}[label=\arabic*.]
  \item $e^{z}=(2{,}7183;\allowbreak\ 1;\allowbreak\ 0{,}3679;\allowbreak\ 7{,}3891)$, suma $11{,}4752$; $p=(0{,}2369;\allowbreak\ 0{,}0871;\allowbreak\ 0{,}0321;\allowbreak\ 0{,}6439)$.
  \item $\ell=-\ln 0{,}6439=0{,}4402$.
  \item $\nabla_z\ell=(0{,}2369;\allowbreak\ 0{,}0871;\allowbreak\ 0{,}0321;\allowbreak\ -0{,}3561)$.
\end{enumerate}
\resultado $p\approx(0{,}2369;\allowbreak\ 0{,}0871;\allowbreak\ 0{,}0321;\allowbreak\ 0{,}6439)$, $\ell\approx 0{,}4402$, $\nabla_z\ell\approx(0{,}2369;\allowbreak\ 0{,}0871;\allowbreak\ 0{,}0321;\allowbreak\ -0{,}3561)$.
\comprobacion el gradiente suma~0 y $\ln 11{,}4752-2=2{,}4402-2=0{,}4402$.

\paragraph{Apartado (d).}
\datos $n=1200$, $k=4$, 5 valores de $\lambda$.
\planteamiento en cada ronda una partición valida y las $k-1$ restantes entrenan; tras elegir $\lambda$ se reentrena con todos los datos (M08, §6.5).
\pasos
\begin{enumerate}[label=\arabic*.]
  \item Cada partición: $1200/4=300$ ejemplos; se entrena con $3\cdot 300=900$.
  \item Entrenamientos: $5\cdot 4=20$, más el final con los $1200$: $21$.
  \item El conjunto de prueba debe quedar intacto: si se usa para elegir $\lambda$, su error deja de ser una estimación insesgada del error de generalización (M08, §1.4).
\end{enumerate}
\resultado \textbf{300 por partición, 900 para entrenar en cada ronda, 21 entrenamientos.}
\comprobacion para cada $\lambda$, cada ejemplo valida exactamente una vez: $4\cdot 300=1200$.

\paragraph{Criterios de corrección.} (a) 1~p $X\T X$ y $X\T y$, 2~p $\hat\theta$, 1~p residuos, $\SSE$ y $\hat\sigma^{2}$. (b) 1~p derivada y paso, 1~p convexidad, 1~p $\alpha^{*}$, 1~p contraejemplo. (c) 1~p $p$, 1~p pérdida y gradiente. (d) 1~p cuentas, 1~p justificación.

\textit{Repasar:} M08, §2.2--2.5; M08, §3.4; M08, §4.3; M08, §6.5.
\end{solucion}

% ===================================================================== P4
\begin{solucion}{eje:ad}
\paragraph{Apartado (a).}
\datos $a=1$, $b=0$, $c=0{,}5$.
\planteamiento traza de evaluación en orden topológico (M10, §1.1).
\pasos
\begin{enumerate}[label=\arabic*.]
  \item Grafo: $a\to d$, $b\to d$, $d\to e$, $a\to f$, $a\to g$, $(e,f,g)\to h$, $(c,h)\to i$, $(f,i)\to j$. Tienen más de una salida $a$ (hacia $d$, $f$ y $g$) y $f$ (hacia $h$ y $j$); sus adjuntos acumulan varias contribuciones.
\begin{center}
\begin{tikzpicture}[font=\small,>={Stealth[length=2mm]},n/.style={draw,circle,minimum size=7mm,inner sep=0pt}]
  \node[n] (a) at (0,1.6) {$a$}; \node[n] (b) at (0,0) {$b$}; \node[n] (c) at (0,-1.6) {$c$};
  \node[n] (d) at (1.8,0.4) {$d$}; \node[n] (e) at (3.4,0.4) {$e$};
  \node[n] (f) at (3.4,2.4) {$f$}; \node[n] (g) at (3.4,1.4) {$g$};
  \node[n] (h) at (5.0,1.0) {$h$}; \node[n] (i) at (6.6,0.0) {$i$}; \node[n] (j) at (8.2,1.2) {$j$};
  \draw[->] (a)--(d); \draw[->] (b)--(d); \draw[->] (d)--(e);
  \draw[->] (a)--(f); \draw[->] (a)--(g);
  \draw[->] (e)--(h); \draw[->] (f)--(h); \draw[->] (g)--(h);
  \draw[->] (h)--(i); \draw[->] (c)--(i); \draw[->] (i)--(j); \draw[->] (f)--(j);
\end{tikzpicture}
\end{center}
  \item $d=1\cdot 0=0$; $e=\sigma(0)=0{,}5$; $f=\ReLU(1)=1$; $g=3$.
  \item $h=2\cdot 0{,}5+1+3=5$; $i=0{,}5\cdot 5=2{,}5$; $j=1+2{,}5^{2}=7{,}25$.
\end{enumerate}
\resultado $d=0,\ e=0{,}5,\ f=1,\ g=3,\ h=5,\ i=2{,}5,\ j=7{,}25$.
\comprobacion $j=f+(c(2\sigma(ab)+\ReLU(a)+3a))^{2}=1+(0{,}5\cdot 5)^{2}=7{,}25$.

\paragraph{Apartado (b).}
\datos valores de (a); $\sigma'(0)=0{,}25$; $\ReLU'(1)=1$.
\planteamiento recorrido inverso con $\bar\jmath=1$; cada nodo suma (\texttt{+=}) las contribuciones de todos sus sucesores (M10, §3.1--3.2): los gradientes se suman en las bifurcaciones~\cite{cs231n}.
\pasos
\begin{enumerate}[label=\arabic*.]
  \item $\bar\imath=\bar\jmath\cdot 2i=5$.
  \item $\bar h=\bar\imath\cdot c=2{,}5$ y $\bar c=\bar\imath\cdot h=25$.
  \item $\bar e=2\bar h=5$; $\bar g=\bar h=2{,}5$; $\bar f=\bar\jmath\cdot 1+\bar h\cdot 1=1+2{,}5=3{,}5$ (dos contribuciones).
  \item $\bar d=\bar e\,\sigma'(d)=5\cdot 0{,}25=1{,}25$; $\bar b=\bar d\cdot a=1{,}25$.
  \item $\bar a=\bar d\cdot b+\bar f\cdot\uno[a\ge 0]+\bar g\cdot 3=0+3{,}5+7{,}5=11$ (tres contribuciones).
  \item Simbólicamente, como en la fuente: $\dfrac{\partial j}{\partial a}=4icb\,\sigma(d)\bigl(1-\sigma(d)\bigr)+(2ic+1)\,\uno[a\ge 0]+6ic$.
\end{enumerate}
\resultado $\partial j/\partial a=11$, $\partial j/\partial b=1{,}25$, $\partial j/\partial c=25$.
\comprobacion la expresión simbólica da $0+(2\cdot 2{,}5\cdot 0{,}5+1)+6\cdot 2{,}5\cdot 0{,}5=3{,}5+7{,}5=11$. Las diferencias centradas con paso $10^{-6}$ dan $11{,}0000$, $1{,}2500$ y $25{,}0000$.

\paragraph{Apartado (c).}
\datos $a=1+\varepsilon$, $b=0$, $c=0{,}5$.
\planteamiento se propaga el par (valor, tangente); la parte en $\varepsilon$ de $j$ es $\partial j/\partial a$ (M10, §2.2).
\pasos
\begin{enumerate}[label=\arabic*.]
  \item $d=(1+\varepsilon)\cdot 0=0$; $e=\sigma(0)=0{,}5$; $f=1+\varepsilon$; $g=3+3\varepsilon$.
  \item $h=1+(1+\varepsilon)+(3+3\varepsilon)=5+4\varepsilon$; $i=2{,}5+2\varepsilon$.
  \item $i^{2}=6{,}25+2\cdot 2{,}5\cdot 2\varepsilon=6{,}25+10\varepsilon$ (pues $\varepsilon^{2}=0$); $j=7{,}25+11\varepsilon$.
  \item El modo directo da una columna del jacobiano por pasada: el gradiente respecto a $(a,b,c)$ exige 3 pasadas, frente a 1 del modo inverso (M10, §2.3 y §4.2).
\end{enumerate}
\resultado \textbf{$j=7{,}25+11\varepsilon$: coincide con el modo inverso. Directo: 3 pasadas; inverso: 1.}
\comprobacion la parte real, $7{,}25$, es el valor de (a).

\paragraph{Apartado (d).}
\datos lote de 128; $E$ ($128\times 30$), $z^{(1)}$ y $h$ ($128\times 64$), $z^{(2)}$ y $p$ ($128\times 28$); $N=4084$; capa oculta $30\to 64$.
\planteamiento contar elementos de los tensores guardados (M10, §5.1); un producto $(B\times n)(n\times k)$ cuesta $2Bnk$ FLOP y la pasada hacia atrás hace dos productos del mismo tamaño (M10, §4.4).
\pasos
\begin{enumerate}[label=\arabic*.]
  \item $E$: $128\cdot 30=3840$; $z^{(1)}$ y $h$: $2\cdot 128\cdot 64=16\,384$; $z^{(2)}$ y $p$: $2\cdot 128\cdot 28=7168$.
  \item Total: $3840+16\,384+7168=27\,392$ números, $27\,392/4084=6{,}71$ veces los parámetros.
  \item Pesos de la capa oculta: $30\cdot 64=1920$. Hacia delante: $2\cdot 128\cdot 1920=491\,520$ FLOP.
  \item Hacia atrás ($\bar EW\T$ y $E\T\bar Z$): $2\cdot 491\,520=983\,040$ FLOP.
  \item Total: $1\,474\,560=6\cdot 1920\cdot 128$.
\end{enumerate}
\resultado \textbf{$27\,392$ activaciones guardadas ($\approx 6{,}7$ veces $N$); $491\,520$ FLOP hacia delante y $983\,040$ hacia atrás, en total $6ND$ con $N=1920$ y $D=128$.}
\comprobacion la proporción $1:2$ entre las dos pasadas es la que da la regla $C\approx 6ND$ (M10, §4.4). Las activaciones crecen con el tamaño del lote y los parámetros no; por eso dominan la memoria del entrenamiento (M10, §5.1).

\paragraph{Criterios de corrección.} (a) 1~p grafo, 1~p valores. (b) 1~p adjuntos hasta $\bar d$, 1~p acumulación en $\bar f$ y $\bar a$, 1~p expresión simbólica. (c) 1~p dual, 1~p número de pasadas. (d) 1,5~p memoria, 1,5~p FLOP.

\textit{Repasar:} M10, §1.1; M10, §2.2--2.3; M10, §3.1--3.2; M10, §4.4; M10, §5.1.
\end{solucion}

% ===================================================================== P5
\begin{solucion}{eje:prof}
\paragraph{Apartado (a).}
\datos $z=(90,88,80)$; máximo de fp32 $3{,}40\cdot 10^{38}=e^{88{,}72}$; clase correcta~1.
\planteamiento $\LSE(z)=c+\ln\sum_ie^{z_i-c}$ con $c=\max_iz_i$ (M11, §6.2) y pérdida $\LSE(z)-z_1$ (M11, §6.3).
\pasos
\begin{enumerate}[label=\arabic*.]
  \item Invariancia: $e^{z_i+c}=e^{c}e^{z_i}$ y el factor $e^{c}>0$ es común al numerador y al denominador, luego se cancela.
  \item Ingenuo: $e^{90}=1{,}22\cdot 10^{39}>3{,}40\cdot 10^{38}$ da \Inf; el denominador también es \Inf\ y $\Inf/\Inf=\NaN$.
  \item $z-90=(0,-2,-10)$; $e^{z-90}=(1;\allowbreak\ 0{,}1353;\allowbreak\ 4{,}540\cdot 10^{-5})$; suma $1{,}1354$; $\ln 1{,}1354=0{,}1270$.
  \item $\LSE(z)=90{,}1270$; $p=(0{,}8808;\allowbreak\ 0{,}1192;\allowbreak\ 4{,}00\cdot 10^{-5})$.
  \item Pérdida: $90{,}1270-88=2{,}1270$.
\end{enumerate}
\resultado \textbf{ingenuo: \Inf\ y \NaN}; $\LSE(z)=90{,}1270$; $p\approx(0{,}8808;\allowbreak\ 0{,}1192;\allowbreak\ 4{,}0\cdot 10^{-5})$; \textbf{pérdida} $2{,}1270$.
\comprobacion $-\ln 0{,}1192=2{,}127$. Cotas: $\max z\le\LSE(z)\le\max z+\ln 3$, es decir, $90\le 90{,}127\le 91{,}099$.

\paragraph{Apartado (b).}
\datos $x=(-1,0,2,3)$, $d=4$, $\epsilon=0$.
\planteamiento $\LN(x)=(x-\mu)/\sigma$ y $\RMSNorm(x)=x/\RMS(x)$ (M11, §3.1); invariancias de M11, §3.2.
\pasos
\begin{enumerate}[label=\arabic*.]
  \item $\mu=1$; $x-\mu=(-2,-1,1,2)$; $\sigma^{2}=10/4=2{,}5$; $\sigma=1{,}5811$; $\LN(x)=(-1{,}2649;\allowbreak\ -0{,}6325;\allowbreak\ 0{,}6325;\allowbreak\ 1{,}2649)$.
  \item $\RMS(x)=\sqrt{14/4}=\sqrt{3{,}5}=1{,}8708$; $\RMSNorm(x)=(-0{,}5345;\allowbreak\ 0;\allowbreak\ 1{,}0690;\allowbreak\ 1{,}6036)$.
  \item $\RMS(10x)=\sqrt{350}=18{,}708=10\RMS(x)$, luego $\RMSNorm(10x)=\RMSNorm(x)$.
  \item LayerNorm también es invariante a la escala: $\mu$ y $\sigma$ se multiplican por~10 y el cociente no cambia.
  \item Traslación: $\LN(x+c\uno)=\LN(x)$, porque la media absorbe $c$; RMSNorm no lo es: $\RMSNorm(x+\uno)=(0;\allowbreak\ 1;\allowbreak\ 3;\allowbreak\ 4)/\sqrt{6{,}5}=(0;\allowbreak\ 0{,}3922;\allowbreak\ 1{,}1767;\allowbreak\ 1{,}5689)$.
\end{enumerate}
\resultado $\LN(x)\approx(-1{,}2649;\allowbreak\ -0{,}6325;\allowbreak\ 0{,}6325;\allowbreak\ 1{,}2649)$, $\RMSNorm(x)\approx(-0{,}5345;\allowbreak\ 0;\allowbreak\ 1{,}0690;\allowbreak\ 1{,}6036)$; \textbf{ambas son invariantes a la escala y solo LN lo es a la traslación.}
\comprobacion la media de los cuadrados de ambas salidas vale~1: $(1{,}6+0{,}4+0{,}4+1{,}6)/4=1$ y $(0{,}2857+0+1{,}1429+2{,}5714)/4=1$.

\paragraph{Apartado (c).}
\datos $g=(2,-1,2)$, $c=1{,}5$; 8 capas con $z=2$.
\planteamiento recorte $g\cdot\min(1,c/\norm{g})$ (M11, §4.2); factor plano $\prod\tanh'(z)$ (M11, §1.3) y residual $(1\pm\tanh'(z))^{8}$ (M11, §2.2).
\pasos
\begin{enumerate}[label=\arabic*.]
  \item $\norm{g}=3>1{,}5$; factor $0{,}5$; $g=(1;\allowbreak\ -0{,}5;\allowbreak\ 1)$.
  \item $\tanh 2=0{,}9640$; $\tanh'(2)=1-0{,}9640^{2}=0{,}07065$.
  \item Sin residuos: $0{,}07065^{8}=6{,}21\cdot 10^{-10}$.
  \item Con residuos: $1{,}07065^{8}=1{,}7266$ y $0{,}92935^{8}=0{,}5565$.
\end{enumerate}
\resultado $g=(1;\allowbreak\ -0{,}5;\allowbreak\ 1)$; \textbf{factor plano} $6{,}21\cdot 10^{-10}$; \textbf{residual} $1{,}7266$ \textbf{($+$) o} $0{,}5565$ \textbf{($-$).}
\comprobacion $\norm{(1;-0{,}5;1)}=\sqrt{2{,}25}=1{,}5=c$. La identidad mantiene el gradiente en orden~1, frente a $10^{-10}$ sin ella.

\paragraph{Apartado (d).}
\datos separación $2^{-8}=0{,}00390625$ en $[0{,}5;\allowbreak\ 1)$; $w=1$; tres actualizaciones independientes; 16 contribuciones de varianza~1.
\planteamiento redondeo al representable más próximo (M11, §5.2); las varianzas de sumandos incorrelados se suman (M11, §2.3).
\pasos
\begin{enumerate}[label=\arabic*.]
  \item Vecinos por debajo de~1: $0{,}99609$ y $0{,}99219$; medio hueco, $2^{-9}=0{,}00195$.
  \item $0{,}999$ dista $0{,}001<0{,}00195$ de~1: se redondea a~1 y la actualización se pierde.
  \item $0{,}997$ dista $0{,}00091$ de $0{,}99609$: se redondea a $0{,}99609$.
  \item $0{,}995$ dista $0{,}00109$ de $0{,}99609$ y $0{,}0028$ de $0{,}99219$: se redondea a $0{,}99609$.
  \item Remedio: copia maestra de los pesos en fp32 y acumulación de las actualizaciones en fp32 (M11, §5.3).
  \item Flujo: $\Var=1+16\cdot 1=17$; desviación típica $\sqrt{17}=4{,}1231$.
\end{enumerate}
\resultado \textbf{$1$; $0{,}9961$; $0{,}9961$ (la primera actualización se pierde y las otras dos dan el mismo peso); desviación típica del flujo $\sqrt{17}\approx 4{,}1231$.}
\comprobacion PyTorch con \texttt{dtype=torch.bfloat16} da $1{,}0$, $0{,}99609375$ y $0{,}99609375$. El crecimiento como $\sqrt L$ motiva la normalización previa y el escalado de las proyecciones residuales (M11, §2.3).

\paragraph{Criterios de corrección.} (a) 1~p invariancia y diagnóstico, 1~p $\LSE$, 1~p $p$ y pérdida. (b) 1~p LN, 1~p RMSNorm, 1~p invariancias. (c) 1~p recorte, 1~p factores. (d) 1~p redondeos y remedio, 1~p desviación típica.

\textit{Repasar:} M11, §6.2--6.3; M11, §3.1--3.2; M11, §4.2; M11, §1.3; M11, §2.2--2.3; M11, §5.2--5.3.
\end{solucion}

% ===================================================================== P6
\begin{solucion}{eje:est}
\paragraph{Apartado (a).}
\datos $B=5$, $\abs{\V}=28$, $d=16$, $H=128$, con sesgos; 800 pasos de 128 ejemplos.
\planteamiento parámetros por capa como en M09, §6.3; $C\approx 6ND$ con $D$ el número de ejemplos procesados (M10, §4.4).
\pasos
\begin{enumerate}[label=\arabic*.]
  \item Embeddings: $28\cdot 16=448$.
  \item Oculta: entrada $5\cdot 16=80$; $80\cdot 128+128=10\,368$.
  \item Salida: $128\cdot 28+28=3612$.
  \item $N=448+10\,368+3612=14\,428$; $D=800\cdot 128=102\,400$; $C=6\cdot 14\,428\cdot 102\,400=8{,}86\cdot 10^{9}$ FLOP.
\end{enumerate}
\resultado $N=14\,428$; $C\approx 8{,}86\cdot 10^{9}$ \textbf{FLOP}, unos segundos en un portátil.
\comprobacion la consulta de embeddings no es un producto de matrices; sin ella, $6\cdot 13\,980\cdot 102\,400=8{,}59\cdot 10^{9}$. La regla es una cota por exceso.

\paragraph{Apartado (b).}
\datos $\abs{\V}=50\,000$, $d=64$, contexto~8, $H=512$.
\planteamiento mismo conteo; AdamW en fp32 guarda pesos, gradientes, $m$ y $v$: $4\cdot 4=16$ bytes por parámetro (M07, §5.3).
\pasos
\begin{enumerate}[label=\arabic*.]
  \item Embeddings: $50\,000\cdot 64=3\,200\,000$.
  \item Oculta: $(8\cdot 64)\cdot 512+512=262\,656$.
  \item Salida: $512\cdot 50\,000+50\,000=25\,650\,000$.
  \item $N=29\,112\,656$; fracción de la salida: $25\,650\,000/29\,112\,656=0{,}8811$.
  \item Memoria: $16\cdot 29\,112\,656=465\,802\,496$ bytes $\approx 466$~MB ($444$~MiB).
\end{enumerate}
\resultado $N=29\,112\,656$; \textbf{la salida es el $88{,}1\,\%$}; $\approx 466$ \textbf{MB}.
\comprobacion embeddings y salida suman el $99{,}1\,\%$: en los modelos pequeños, el vocabulario domina el recuento (M08, §4.5).

\paragraph{Apartado (c).}
\datos $N=2{,}911\cdot 10^{7}$, $D=10^{8}$, $10^{12}$~FLOP/s; $7\cdot 10^{9}$ parámetros a 16 bytes; 80~GB.
\planteamiento tiempo $=6ND/(\text{FLOP/s})$; memoria $=16N$ bytes.
\pasos
\begin{enumerate}[label=\arabic*.]
  \item $C=6\cdot 2{,}911\cdot 10^{7}\cdot 10^{8}=1{,}747\cdot 10^{16}$ FLOP.
  \item Tiempo: $1{,}747\cdot 10^{16}/10^{12}=17\,468$~s $=4{,}85$~h.
  \item $16\cdot 7\cdot 10^{9}=1{,}12\cdot 10^{11}$ bytes $=112$~GB $>80$~GB: no cabe.
\end{enumerate}
\resultado $\approx 1{,}75\cdot 10^{16}$ \textbf{FLOP, unas 4,9~h; el modelo de 7000 millones necesita 112~GB y no cabe en 80~GB}. Habría que repartir el estado del optimizador entre varias GPU o reducir la precisión (anticipo del módulo~17).
\comprobacion orden de magnitud: $6\cdot 3\cdot 10^{7}\cdot 10^{8}=1{,}8\cdot 10^{16}$ y $1{,}8\cdot 10^{4}$~s $\approx 5$~h.

\paragraph{Criterios de corrección.} (a) 2~p conteo, 1~p FLOP. (b) 2~p conteo, 1~p fracción, 1~p memoria. (c) 2~p tiempo, 1~p conclusión. Se aceptan resultados dentro de un factor~2 si el razonamiento es correcto.

\textit{Repasar:} M09, §6.3; M10, §4.4; M07, §5.3; M08, §4.5.
\end{solucion}
